% !TeX root = ../main.tex
\chapter{Rock Mass Classification}
\label{ch:classification}

Rock mass classification occupies an intermediate position in the design
sequence. It converts observations that can be recorded in the field into
ratings that empirical design methods accept as input, and it does so through
schemes calibrated against case histories rather than derived from mechanics.
Three systems dominate practice: the Rock Mass Rating of
\citet{bieniawski1989rmr}, the $Q$-system of
\citet{barton1974q}, and the Geological Strength Index introduced
by Hoek and colleagues for parameter estimation rather than for support
selection~\citep{marinos2000gsi}. This chapter sets out the three, the data
each requires, and the reasons their ratings diverge. The support
recommendations built on the ratings, and the conditions under which those
recommendations may legitimately be compared, are taken up in
Chapter~\ref{ch:empirical-design}.

\section{Classification systems}
\label{sec:classification-systems}

\subsection{Purpose and workflow of classification}
\label{sec:purpose-and-workflow-of-classification}

Classification systems were developed to do five things: to identify the
parameters that most influence rock mass behaviour; to divide rock masses into
groups of similar engineering behaviour; to supply quantitative guidance for
preliminary design; to give engineers and geologists a common vocabulary; and
to provide input to empirical support methods and to rock mass failure
criteria. The line of development runs from the rock load factor of
\citet{terzaghi1946rockload}, through the Rock Mass Rating and the $Q$-system
in the 1970s, to the Geological Strength Index two decades
later~\citep{hoek1995support}. The earlier schemes were written to select
support. The last was written to feed a failure criterion, and the difference
in purpose explains most of the differences in what each system counts.

A classification is not a calculation alone. Each system reduces to arithmetic
once its ratings are assigned, but the ratings themselves are geological
judgements, and a reported RMR, $Q$ or GSI value is incomplete without the
justification for each rating behind it.

Figure~\ref{fig:classification_workflow} sets out the sequence. The data are
collected once and the three systems are evaluated from them in turn; the last
step compares the results, and a comparison that exposes an inconsistency
returns the process to the data rather than to the arithmetic.

\begin{figure}[htbp]
\centering
\resizebox{\textwidth}{!}{%
\begin{tikzpicture}[node distance=0.4cm, font=\small]

\node[smallworkflowbox, wfinput] (data)
{\textbf{1. Collect data}\\
\footnotesize UCS, RQD, joints, water, stress};

\node[smallworkflowbox, wfstep, right=of data] (rmr)
{\textbf{2. Calculate RMR}\\
\footnotesize \(R_1+R_2+R_3+R_4+R_5\)};

\node[smallworkflowbox, wfstep, right=of rmr] (q)
{\textbf{3. Calculate Q}\\
\footnotesize \(RQD/J_n \times J_r/J_a \times J_w/SRF\)};

\node[smallworkflowbox, wfstep, right=of q] (gsi)
{\textbf{4. Estimate GSI}\\
\footnotesize chart + correlations};

\node[smallworkflowbox, wfadopt, right=of gsi] (check)
{\textbf{5. Compare}\\
\footnotesize justify final design value};

\draw[workflowarrow] (data) -- (rmr);
\draw[workflowarrow] (rmr) -- (q);
\draw[workflowarrow] (q) -- (gsi);
\draw[workflowarrow] (gsi) -- (check);

\draw[feedbackarrow] (check.south) -- ++(0,-0.6) -| node[pos=0.25,below,font=\footnotesize]
{revise an input the data contradict, e.g.\ RQD against measured spacing} (data.south);

\end{tikzpicture}
}
\caption{Recommended rock mass classification workflow.}
\label{fig:classification_workflow}
\end{figure}

The three systems measure related but not identical aspects of rock mass
quality, and they should not be forced to agree. A difference between them is
not necessarily a mistake: RMR and $Q$ both penalise groundwater, for example,
while GSI by construction does not. A sound classification explains why the
values differ and then selects a defensible design value.

\subsection{Rock Mass Rating (RMR) system}
\label{sec:rmr-system}

The Rock Mass Rating of \citet{bieniawski1989rmr} is the sum of five ratings,
\[
\text{RMR} = R_1 + R_2 + R_3 + R_4 + R_5 ,
\]
for intact rock strength, RQD, discontinuity spacing, discontinuity condition
and groundwater respectively. The sum of the five as assessed is the
\emph{basic} RMR. A sixth term, an adjustment for the orientation of the
discontinuities relative to the excavation, is applied afterwards for design
applications and is treated with them in Chapter~\ref{ch:empirical-design}.
Neither of these is the quantity from which GSI is estimated; that is a third
form of the rating, defined in Section~\ref{sec:geological-strength-index}.
Adjacent bands in the rating tables that follow share their limits; a value
lying exactly on a shared limit is assigned to the band of which it is the
lower limit, so that an RQD of 90\,\% rates 20
(Figure~\ref{fig:rmr-ratings}).

\begin{figure}[htbp]
\centering
\datafigure[\textwidth]{ch10/rmr-ratings}
\caption{The four numerical RMR ratings as step functions of Tables~\ref{tab:rmr_strength} to~\ref{tab:rmr_groundwater}, with the four rocks of Table~\ref{tab:classification-input-data} where the input is known. $R_5 = 15$ applies only to $u/\sigma_1 = 0$; the limestone's 0.16 rates 7.}
\label{fig:rmr-ratings}
\end{figure}

\paragraph{R1: strength of intact rock material}
The first rating is read from the uniaxial compressive strength or, where that
has not been measured, from the point load index
(Table~\ref{tab:rmr_strength}). The two columns of the table are paired at a
ratio of 25, whereas the ratio measured on a given rock varies with lithology.
Where both are available the uniaxial compressive strength governs: the point
load index is a field substitute, and a disagreement of one rating band
between the two is to be expected.

\begin{table}[htbp]
\caption{RMR rating for intact rock strength.}
\label{tab:rmr_strength}
\centering
\small
\begin{tabular}{l l c}
\toprule
\textbf{Point Load \(I_{s(50)}\) (MPa)} & \textbf{UCS (MPa)} & \textbf{Rating} \\
\midrule
\(>10\) & \(>250\) & 15 \\
4--10 & 100--250 & 12 \\
2--4 & 50--100 & 7 \\
1--2 & 25--50 & 4 \\
-- & 5--25 & 2 \\
-- & 1--5 & 1 \\
-- & \(<1\) & 0 \\
\bottomrule
\end{tabular}
\end{table}

\paragraph{R2: Rock Quality Designation (RQD)}
The second rating is read from the RQD (Table~\ref{tab:rmr_rqd}), defined in
Chapter~\ref{ch:site-investigation} as
\[
\text{RQD} =
\frac{
\sum \text{length of sound core pieces} > 100~\text{mm}
}
{\text{total core run length}}
\times 100\% .
\]

\begin{table}[htbp]
\caption{RMR rating for RQD.}
\label{tab:rmr_rqd}
\centering
\small
\begin{tabular}{c c l}
\toprule
\textbf{RQD (\%)} & \textbf{Rating} & \textbf{Description} \\
\midrule
90--100 & 20 & Excellent \\
75--90 & 17 & Good \\
50--75 & 13 & Fair \\
25--50 & 8 & Poor \\
\(<25\) & 3 & Very poor \\
\bottomrule
\end{tabular}
\end{table}

RQD is a geometric index and not a strength. It records the proportion of core
recovered in pieces longer than 100\,mm, and it says nothing directly about
joint roughness, infilling, shear strength or deformability, which is why the
system carries a separate rating for discontinuity condition.

\paragraph{R3: spacing of discontinuities}
The third rating is read from the spacing of the discontinuities
(Table~\ref{tab:rmr_spacing}).

\begin{table}[htbp]
\caption{RMR rating for discontinuity spacing.}
\label{tab:rmr_spacing}
\centering
\small
\begin{tabular}{c c}
\toprule
\textbf{Joint spacing (m)} & \textbf{Rating} \\
\midrule
\(>2\) & 20 \\
0.6--2 & 15 \\
0.2--0.6 & 10 \\
0.06--0.2 & 8 \\
\(<0.06\) & 5 \\
\bottomrule
\end{tabular}
\end{table}

Where spacing has not been measured directly it may be estimated from RQD
through the linear approximation of \citet{priest1976discontinuity},
\[
RQD = 110.4 - 3.68\lambda
\qquad
(6 < \lambda < 16) ,
\]
where \(\lambda\) is the discontinuity frequency in joints per metre, the mean
spacing being \(s = 1/\lambda\). An RQD of 68\,\% gives
\[
\lambda = \frac{110.4 - 68}{3.68} = 11.52~\text{joints/m},
\qquad
s = \frac{1}{11.52}=0.087~\text{m},
\]
which falls in the 0.06--0.2\,m band of Table~\ref{tab:rmr_spacing}. The
approximation holds only for \(6 < \lambda < 16\), that is for RQD between
about 51\,\% and 88\,\% (Figure~\ref{fig:rqd-frequency}). Outside that range,
and in very poor rock in particular, it gives no more than a rough guide, and
a spacing obtained from it carries that uncertainty into the rating.

\begin{figure}[htbp]
\centering
\datafigure[0.62\textwidth]{ch10/rqd-frequency}
\caption{RQD against discontinuity frequency: the exponential relation of Priest and Hudson for randomly spaced joints, and its linear approximation, valid over the shaded range $6 < \lambda < 16$. The point is the example RQD of 68\,\%, $\lambda = 11.5$~m$^{-1}$.}
\label{fig:rqd-frequency}
\end{figure}

\paragraph{R4: condition of discontinuities}
The fourth rating describes the surfaces themselves
(Table~\ref{tab:rmr_condition}). Field descriptions seldom match a single row.
Where a description mixes features, the rating adopted is the one that best
represents the likely engineering behaviour, with the reason recorded; the
most favourable reading of a mixed description is not a defensible default.

\begin{table}[htbp]
\caption{RMR rating for discontinuity condition.}
\label{tab:rmr_condition}
\centering
\small
\begin{tabularx}{\textwidth}{Y c}
\toprule
\textbf{Condition description} & \textbf{Rating} \\
\midrule
Very rough surfaces, not continuous, no separation, unweathered wall rock & 30 \\
Slightly rough surfaces, separation \(<1\) mm, slightly weathered walls & 25 \\
Slightly rough surfaces, separation \(<1\) mm, highly weathered walls & 20 \\
Slickensided surfaces OR gouge \(<5\) mm thick OR separation 1--5 mm, continuous & 10 \\
Soft gouge \(>5\) mm thick OR separation \(>5\) mm, continuous & 0 \\
\bottomrule
\end{tabularx}
\end{table}

\paragraph{R5: groundwater conditions}
The fifth rating is read from the inflow, from the ratio of joint water
pressure to major principal stress, or from a general description of the
condition (Table~\ref{tab:rmr_groundwater}).

\begin{table}[htbp]
\caption{RMR rating for groundwater.}
\label{tab:rmr_groundwater}
\centering
\small
\begin{tabular}{l l l c}
\toprule
\textbf{Inflow per 10 m tunnel (L/min)} & \textbf{\(u/\sigma_1\)} & \textbf{General condition} & \textbf{Rating} \\
\midrule
None & 0 & Completely dry & 15 \\
\(<10\) & 0--0.1 & Damp & 10 \\
10--25 & 0.1--0.2 & Wet & 7 \\
25--125 & 0.2--0.5 & Dripping & 4 \\
\(>125\) & \(>0.5\) & Flowing & 0 \\
\bottomrule
\end{tabular}
\end{table}

\paragraph{Rock mass classes based on RMR}
The total places the rock mass in one of five classes
(Table~\ref{tab:rmr_classes}). The cohesion and friction angle attached to
each class are indicative values for the rock mass as a whole, obtained by
back-analysis for broad classes at particular scales and confinements. They
pre-date the parameter derivation of Chapter~\ref{ch:rock-mass-parameters} and
are no substitute for it: design values of $c$ and $\phi$ are obtained there by
fitting the Hoek--Brown envelope over the confining stress range of the
problem.

\begin{table}[htbp]
\caption{Rock mass classes based on RMR.}
\label{tab:rmr_classes}
\centering
\small
\begin{tabular}{c c l c c}
\toprule
\textbf{RMR} & \textbf{Class} & \textbf{Description} & \textbf{Cohesion (kPa)} & \textbf{Friction angle} \\
\midrule
81--100 & I & Excellent & \(>400\) & \(>45^\circ\) \\
61--80 & II & Good & 300--400 & 35--45\(^{\circ}\) \\
41--60 & III & Fair & 200--300 & 25--35\(^{\circ}\) \\
21--40 & IV & Poor & 100--200 & 15--25\(^{\circ}\) \\
\(<21\) & V & Very poor & \(<100\) & \(<15^\circ\) \\
\bottomrule
\end{tabular}
\end{table}

\subsection{NGI Q-system}
\label{sec:q-system}

The $Q$-system of \citet{barton1974q} is a product of three quotients,
\[
Q =
\frac{RQD}{J_n}
\times
\frac{J_r}{J_a}
\times
\frac{J_w}{SRF} ,
\]
each of which has a physical reading and its own range
(Figure~\ref{fig:q-quotient-ranges}). The first, $RQD/J_n$, is a measure of
block size. The second, $J_r/J_a$, is a measure of the shear strength between
blocks. The third, $J_w/SRF$, describes the active stress, in which water
pressure and the stress state both act. RQD enters as a percentage: an RQD of
70\,\% is entered as 70, not 0.70, and the decimal form returns a $Q$ value
one hundred times too small.

\begin{figure}[htbp]
\centering
\datafigure[\textwidth]{ch10/q-quotient-ranges}
\caption{Range of each quotient of the $Q$-system between the extreme entries of Tables~\ref{tab:jn} to~\ref{tab:srf}, with RQD from 10 to 100; the number is the span in orders of magnitude. $J_w/SRF$ spans the widest range, which is why the choice of SRF category dominates. Grey lines mark the class limits of Table~\ref{tab:q_quality}.}
\label{fig:q-quotient-ranges}
\end{figure}

\paragraph{Joint set number \(J_n\)}
The joint set number counts the sets present, with an allowance for random
joints (Table~\ref{tab:jn}).

\begin{table}[htbp]
\caption{Joint set number \(J_n\).}
\label{tab:jn}
\centering
\small
\begin{tabularx}{0.85\textwidth}{Y c}
\toprule
\textbf{Joint sets} & \textbf{\(J_n\)} \\
\midrule
Massive, no or few joints & 0.5--1 \\
One joint set & 2 \\
One joint set plus random joints & 3 \\
Two joint sets & 4 \\
Two joint sets plus random joints & 6 \\
Three joint sets & 9 \\
Three joint sets plus random joints & 12 \\
Four or more joint sets, heavily jointed & 15 \\
Crushed rock, earthlike & 20 \\
\bottomrule
\end{tabularx}
\end{table}

\paragraph{Joint roughness number \(J_r\)}
The joint roughness number describes the surface at both the small scale and
the scale of the undulation (Table~\ref{tab:jr}).

\begin{table}[htbp]
\caption{Joint roughness number \(J_r\).}
\label{tab:jr}
\centering
\small
\begin{tabularx}{0.85\textwidth}{Y c}
\toprule
\textbf{Description} & \textbf{\(J_r\)} \\
\midrule
Discontinuous joints & 4 \\
Rough or irregular, undulating & 3 \\
Smooth, undulating & 2 \\
Slickensided, undulating & 1.5 \\
Rough or irregular, planar & 1.5 \\
Smooth, planar & 1 \\
Slickensided, planar & 0.5 \\
\bottomrule
\end{tabularx}
\end{table}

\paragraph{Joint alteration number \(J_a\)}
The joint alteration number describes the filling and the degree of rock-wall
contact, in three groups according to whether the walls remain in contact
during shear (Table~\ref{tab:ja}). The table is abridged from the original.

\begin{table}[htbp]
\caption{Joint alteration number \(J_a\).}
\label{tab:ja}
\centering
\small
\begin{tabularx}{\textwidth}{Y c c}
\toprule
\textbf{Description} & \textbf{Residual friction angle} & \textbf{\(J_a\)} \\
\midrule
\multicolumn{3}{l}{\textbf{(a) Rock-wall contact, no mineral fillings, only coatings}} \\
Tightly healed, hard, non-softening, impermeable filling & -- & 0.75 \\
Unaltered joint walls, surface staining only & 25--35\(^{\circ}\) & 1 \\
Slightly altered, non-softening mineral coating, sandy particles & 25--30\(^{\circ}\) & 2 \\
Silty- or sandy-clay coatings, small clay fraction, non-softening & 20--25\(^{\circ}\) & 3 \\
Softening or low-friction mineral coatings & 8--16\(^{\circ}\) & 4 \\
\midrule
\multicolumn{3}{l}{\textbf{(b) Rock-wall contact before 10 cm shear, thin fillings}} \\
Sandy particles, clay-free disintegrated rock & 25--30\(^{\circ}\) & 4 \\
Strongly over-consolidated, non-softening clay fillings & 16--24\(^{\circ}\) & 6 \\
Medium or low over-consolidation, softening clay fillings & 12--16\(^{\circ}\) & 8 \\
Swelling clay fillings & 6--12\(^{\circ}\) & 8--12 \\
\midrule
\multicolumn{3}{l}{\textbf{(c) No rock-wall contact when sheared, thick fillings}} \\
Zones of disintegrated or crushed rock, strongly over-consolidated & 16--24\(^{\circ}\) & 6 \\
Zones of clay, medium or low over-consolidation & 12--16\(^{\circ}\) & 8 \\
Zones of clay, swelling clay & 6--12\(^{\circ}\) & 8--12 \\
Thick continuous zones of clay, strongly over-consolidated & 16--24\(^{\circ}\) & 10 \\
Thick continuous zones of clay, medium to low over-consolidation & 12--16\(^{\circ}\) & 13 \\
Thick continuous zones with swelling clay & 6--12\(^{\circ}\) & 13--20 \\
\bottomrule
\end{tabularx}
\end{table}

\paragraph{Joint water reduction factor \(J_w\)}
The joint water reduction factor is read from the inflow and the approximate
water pressure (Table~\ref{tab:jw}).

\begin{table}[htbp]
\caption{Joint water reduction factor \(J_w\). The approximate water
pressures, given in the original as 1, 2.5 and 10~kgf/cm\(^2\), are converted
at 1~kgf/cm\(^2\) = 0.098~MPa and rounded.}
\label{tab:jw}
\centering
\small
\begin{tabularx}{\textwidth}{c Y c c}
\toprule
\textbf{Code} & \textbf{Conditions} & \textbf{Pressure (MPa)} & \textbf{\(J_w\)} \\
\midrule
A & Dry excavations or minor inflow & \(<0.1\) & 1.0 \\
B & Medium inflow, occasional outwash of joint fillings & 0.1--0.25 & 0.66 \\
C & Jet inflow or high pressure in competent rock & 0.25--1.0 & 0.5 \\
D & Large inflow or high pressure, considerable outwash & 0.25--1.0 & 0.33 \\
E & Exceptionally high inflow, decaying with time & \(>1.0\) & 0.2--0.1 \\
F & Exceptionally high inflow, continuing without decay & \(>1.0\) & 0.1--0.05 \\
\bottomrule
\end{tabularx}
\end{table}

\paragraph{Stress reduction factor (SRF)}
\label{sec:srf}
The stress reduction factor accounts for weakness zones, stress problems,
squeezing and swelling (Table~\ref{tab:srf}), and it is among the most
judgement-sensitive parts of the system.

\begin{table}[htbp]
\caption{Stress Reduction Factor (SRF).}
\label{tab:srf}
\centering
\small
\begin{tabularx}{\textwidth}{Y c}
\toprule
\textbf{Description} & \textbf{SRF} \\
\midrule
\multicolumn{2}{l}{\textbf{(a) Weakness zones intersecting the excavation}} \\
Multiple weakness zones with clay, very loose surrounding rock & 10 \\
Multiple shear zones in competent clay-free rock & 7.5 \\
Single weak zone with or without clay, depth \(\leq 50\) m & 5.0 \\
Loose, open joints, heavily jointed or ``sugar cube'' rock mass & 5.0 \\
Single weak zone with or without clay, depth \(>50\) m & 2.5 \\
\midrule
\multicolumn{2}{l}{\textbf{(b) Competent rock, stress problems}} \\
Low stress, near surface, open joints, \(\sigma_{ci}/\sigma_1^{IS} > 200\) & 2.5 \\
Medium stress, favourable conditions, \(\sigma_{ci}/\sigma_1^{IS}=200\)--10 & 1.0 \\
High stress, very tight structure, \(\sigma_{ci}/\sigma_1^{IS}=10\)--5 & 0.5--2 \\
Moderate spalling/slabbing, \(\sigma_{ci}/\sigma_1^{IS}=5\)--3 & 5--50 \\
Spalling or rock burst, \(\sigma_{ci}/\sigma_1^{IS}=3\)--2 & 50--200 \\
Heavy rock burst, \(\sigma_{ci}/\sigma_1^{IS}<2\) & 200--400 \\
\midrule
\multicolumn{2}{l}{\textbf{(c) Squeezing rock}} \\
Mild squeezing, \(\sigma_{\theta}/\sigma_c = 1\)--5 & 5--10 \\
Heavy squeezing, \(\sigma_{\theta}/\sigma_c > 5\) & 10--20 \\
\midrule
\multicolumn{2}{l}{\textbf{(d) Swelling rock}} \\
Mild swelling & 5--10 \\
Heavy swelling & 10--15 \\
\bottomrule
\end{tabularx}
\end{table}

The choice of category matters more than the value read within it. Category
(a) applies wherever the excavation intersects weakness zones, heavily jointed
zones or poorly interlocked rock, whatever the stress level. Category (b)
applies when the rock mass is competent and the ratio
\(\sigma_{ci}/\sigma_1^{IS}\) can be meaningfully calculated. In that ratio
\(\sigma_1^{IS}\) is the major in-situ principal stress, the demand, and not
the major principal stress at failure \(\sigma_1^f\) returned by a failure
criterion, the capacity. Chapter~\ref{ch:rock-mechanics} keeps the two apart
by the same superscripts, and substituting one for the other here returns the
wrong stress reduction factor.

\paragraph{Rock mass quality based on Q}
The resulting value places the rock mass on a logarithmic scale of quality
(Table~\ref{tab:q_quality}).

\begin{table}[htbp]
\caption{Rock mass quality based on Q-value.}
\label{tab:q_quality}
\centering
\small
\begin{tabular}{c l}
\toprule
\textbf{Q-value} & \textbf{Rock mass quality} \\
\midrule
\(>400\) & Exceptionally good \\
100--400 & Extremely good \\
40--100 & Very good \\
10--40 & Good \\
4--10 & Fair \\
1--4 & Poor \\
0.1--1 & Very poor \\
0.01--0.1 & Extremely poor \\
\(<0.01\) & Exceptionally poor \\
\bottomrule
\end{tabular}
\end{table}

\section{Geological Strength Index}
\label{sec:geological-strength-index}

The Geological Strength Index describes two things only: the structure of the
rock mass and the condition of its discontinuity surfaces. It is the input
through which geology enters the generalised Hoek--Brown
criterion~\citep{hoek2019hb} of Chapter~\ref{ch:rock-mass-parameters}, and for
that reason it deliberately excludes groundwater, the orientation of the
excavation and temporary construction effects, each of which is accounted for
elsewhere in an analysis that uses the criterion.

\paragraph{Structure and surface condition}
Structure is described on a scale of decreasing interlock. An \emph{intact or
massive} rock mass has few, widely spaced discontinuities. A \emph{blocky}
mass is well interlocked, with cubical blocks formed by about three joint
sets; a \emph{very blocky} mass is interlocked but partially disturbed, with
multifaceted angular blocks. A \emph{blocky, disturbed or seamy} mass is
folded or faulted, with many intersecting discontinuities. A
\emph{disintegrated} mass is a poorly interlocked, heavily broken mixture of
fragments, and a \emph{laminated or sheared} mass lacks blockiness altogether
because of weak foliation, schistosity or shearing.

Surface condition is described on a second scale. \emph{Very good} surfaces
are very rough, fresh and unweathered. \emph{Good} surfaces are rough,
slightly weathered and iron-stained. \emph{Fair} surfaces are smooth,
moderately weathered and altered. \emph{Poor} surfaces are slickensided and
highly weathered, with compact coatings or fillings, and \emph{very poor}
surfaces are slickensided and highly weathered with soft clay coatings or
fillings.

\paragraph{Chart-based GSI estimation}
The two descriptions are combined on the GSI chart, reproduced numerically as
Table~\ref{tab:gsi} of Appendix~\ref{app:constants}, to give a range rather
than a number. The range is the result. A single GSI value reported without
its descriptive basis cannot be checked, and a record of the form ``blocky to
very blocky, with slightly rough but highly weathered surfaces, giving a GSI
range from which a representative value is adopted'' is the minimum that
allows another assessor to agree or to disagree.

\paragraph{GSI from RMR}
Where a chart assessment is unavailable or needs checking, GSI may be
estimated from the Rock Mass Rating as
\[
\text{GSI} = RMR'_{89} - 5 ,
\]
in which \(RMR'_{89}\) is the 1989 rating computed with the groundwater rating
set to 15, the value for completely dry conditions, and with no orientation
adjustment. It is related to the basic rating, which carries the groundwater
rating \(R_5\) actually assessed, by
\[
RMR'_{89} = RMR_{\text{basic}} - R_5 + 15 .
\]
For a basic RMR of 60 assessed with \(R_5 = 7\),
\[
RMR'_{89} = 60 - 7 + 15 = 68,
\qquad
\text{GSI} = 68 - 5 = 63 .
\]
Entering the basic rating directly would give 55. The eight-point difference
is the groundwater penalty counted twice: once in the classification, and
again when pore pressure is introduced explicitly in the analysis for which
the GSI is wanted.

\paragraph{Other empirical correlations}
Two further correlations estimate GSI from quantities already assessed
(Figure~\ref{fig:gsi-correlations}),
\[
GSI = 1.5J_{\text{cond89}} + \frac{RQD}{2}
\]
and
\[
GSI =
\frac{52(J_r/J_a)}
{1+(J_r/J_a)}
+
\frac{RQD}{2} ,
\]
where \(J_{\text{cond89}}\) is the RMR discontinuity condition rating \(R_4\).
It is not the joint roughness number \(J_r\) of the $Q$-system, which appears
in the second expression, and the similarity of the two symbols is a known
source of error.

\begin{figure}[htbp]
\centering
\datafigure[\textwidth]{ch10/gsi-correlations}
\caption{GSI from the two correlations, for RQD of 25, 50, 75 and 100\,\%: from the RMR condition rating $J_{\mathrm{cond89}}$ (left) and from the $Q$-system quotient $J_r/J_a$ (right). The square is the sandstone of Table~\ref{tab:classification-input-data}, $J_{\mathrm{cond89}} = 20$ and RQD 68\,\%, giving 64 against 63 from $RMR'_{89}$.}
\label{fig:gsi-correlations}
\end{figure}

\paragraph{Selecting a GSI value}
Four estimates are therefore available: the chart value, the value from
\(RMR'_{89}\), and the two from the correlations above.
Figure~\ref{fig:gsi_selection} shows how they are reconciled. The chart-based
value is the primary one, because it is the direct geological judgement the
index was defined to record, and the correlations serve as checks on it. Where
the estimates differ significantly, the difference is itself information about
the rock mass or about the logging, and it is explained before a cautious
value within the supported range is adopted.

\begin{figure}[htbp]
\centering
\resizebox{0.95\textwidth}{!}{%
\begin{tikzpicture}[node distance=0.45cm, font=\small]

\node[workflowbox, wfstep] (chart)
{\textbf{Chart GSI}\\
\footnotesize structure + surface condition};

\node[workflowbox, wfstep, right=of chart] (corr)
{\textbf{Empirical GSI}\\
\footnotesize RMR, RQD, \(J_r/J_a\)};

\node[workflowbox, wfcheck, right=of corr] (compare)
{\textbf{Compare}\\
\footnotesize check consistency and outliers};

\node[workflowbox, wfadopt, right=of compare] (select)
{\textbf{Adopt final GSI}\\
\footnotesize justify design value or range};

\draw[workflowarrow] (chart) -- (corr);
\draw[workflowarrow] (corr) -- (compare);
\draw[workflowarrow] (compare) -- (select);

\end{tikzpicture}
}
\caption{Workflow for selecting a defensible GSI value.}
\label{fig:gsi_selection}
\end{figure}

\section{An illustrative data set}
\label{sec:classification-design-data}

Table~\ref{tab:classification-input-data} sets out the information the three
systems require for four rock masses spanning the range of civil practice: a
strong granite at depth under high horizontal stress, a sandstone and a
mudstone at intermediate and shallow depth, and a limestone near surface. The
same four are carried through the stress and strength calculations of
Chapter~\ref{ch:rock-mechanics} and the parameter derivation of
Chapter~\ref{ch:rock-mass-parameters}.

\begin{table}[htbp]
\caption{Rock mass information required by the three classification systems.}
\label{tab:classification-input-data}
\centering
\scriptsize
\begin{tabularx}{\textwidth}{l Y Y Y Y}
\toprule
\textbf{Parameter} & \textbf{Granite} & \textbf{Sandstone} & \textbf{Mudstone} & \textbf{Limestone} \\
\midrule
\(\sigma_{ci}\) (MPa) & 265 & 110 & 27 & 85 \\
\(I_{s(50)}\) (MPa) & 12.2 & 5.3 & 2.3 & 5.0 \\
RQD (\%) & 70 & 68 & 38 & 52 \\
Joint condition & Tight, closed, very rough, undulating, unweathered & Slightly rough, highly weathered, not separated & Slickensided and planar, aperture 1--2 mm & Smooth, undulating, slightly weathered, clay-filled, aperture 1--2 mm \\
Joint sets / structure & Two joint sets from core observations & One sub-horizontal and one sub-vertical set & Poorly interlocked, heavily jointed, more than four sets & One joint set, spacing approximately 0.06 m \\
Groundwater & Dry & Medium inflow expected & Dripping at outcrop; medium inflow expected & Groundwater table 10 m below surface \\
Depth (m) & 780 & 450 & 100 & 80 \\
Stress ratio \(k\) & 3.2 & 1.8 & 1.0 & 2.0 \\
\bottomrule
\end{tabularx}
\end{table}

The groundwater entries are recorded in the language of the field and must be
translated into the quantities the two systems use. ``Dry'' is the completely
dry condition of Table~\ref{tab:rmr_groundwater} (\(R_5 = 15\)) and category~A
of Table~\ref{tab:jw} (\(J_w = 1.0\)). ``Medium inflow'' is category~B
(\(J_w = 0.66\)) but has no counterpart in the RMR table; for the sandstone it
is read as the wet condition, 10--25~L/min per 10~m, giving \(R_5 = 7\) and a
basic RMR of 60 (Figure~\ref{fig:rmr-sandstone}). For
the mudstone the observed dripping governs (\(R_5 = 4\)), with \(J_w = 0.66\)
for the medium inflow. The limestone entry needs a conversion. At 80~m depth
with the water table 10~m below surface the head is 70~m, so
\(u = 70 \times 9.81 = 687\)~kPa, or 0.69~MPa, and with
\(\sigma_1 = k\gamma z = 2.0 \times 0.027 \times 80 = 4.32\)~MPa
(\(\gamma\) from Table~\ref{tab:intact_project_data}),
\(u/\sigma_1 = 0.69/4.32 = 0.16\) and \(R_5 = 7\). The ratio is taken to the
major principal stress, as the table defines it, and not to the vertical
stress, which here is the minor one; where \(k\) departs from unity the
difference changes the rating. For the $Q$-system, if the excavation does not
drain the joints, 0.69~MPa lies in the 0.25--1.0~MPa band of categories~C
and~D (\(J_w = 0.5\) or 0.33), the choice between them
resting on the inflow and outwash observed.

\begin{figure}[htbp]
\centering
\datafigure[0.62\textwidth]{ch10/rmr-sandstone}
\caption{The five RMR ratings of the sandstone: $R_1 = 12$ ($\sigma_{ci} = 110$~MPa), $R_2 = 13$ (RQD 68\,\%), $R_3 = 8$ ($s = 0.087$~m), $R_4 = 20$ and $R_5 = 7$. Replacing $R_5$ by the dry value of 15 gives $RMR'_{89} = 68$ and GSI 63; the basic rating would give 55.}
\label{fig:rmr-sandstone}
\end{figure}

The point load indices of the mudstone and the limestone fall one rating band
above their uniaxial compressive strengths. That is the disagreement
anticipated in Section~\ref{sec:rmr-system}, and the uniaxial strength
governs the rating.

A summary of this kind cannot carry everything a rating depends on. Ratings
assigned from tabulated descriptors inherit whatever the descriptors omit, and
the justification for each rating, with the calculation behind each derived
value, belongs in the record alongside the result.

\chapterreferences
